% -----------------------------------------------------------------------------
% DEXPI Beamer Theme
% -----------------------------------------------------------------------------
% Copyright (c) 2026 DEXPI e.V.
%
% This work is licensed under the Creative Commons Attribution-ShareAlike 4.0
% International License (CC BY-SA 4.0).
%
% You are free to:
%   - Share: copy and redistribute the material in any medium or format
%   - Adapt: remix, transform, and build upon the material for any purpose
%
% Under the following terms:
%   - Attribution: You must give appropriate credit to DEXPI e.V.,
%     provide a link to the license, and indicate if changes were made.
%   - ShareAlike: If you remix, transform, or build upon the material,
%     you must distribute your contributions under the same license.
%
% Full license text:
%   https://creativecommons.org/licenses/by-sa/4.0/
%
%  Author: Michael Wiedau
% -----------------------------------------------------------------------------

%----------------------------------------------------------------------------------------
%    PACKAGES AND THEMES
%----------------------------------------------------------------------------------------

\documentclass[aspectratio=169,xcolor=dvipsnames]{beamer}
\usetheme{DEXPI}

\usepackage{hyperref}
\usepackage{graphicx}
\usepackage{booktabs}
\usepackage{pgfgantt}
\usepackage[portuguese]{babel}


\usepackage{tikz}
\usepackage{circuitikz}
\usetikzlibrary{shapes.geometric,calc,positioning,decorations.pathreplacing}
\usepackage{pgfplots}
\pgfplotsset{compat=1.10}
\usepackage{color}

\newcommand{\hrefcol}[2]{\textcolor{cyan}{\href{#1}{#2}}}



% adicionar o numero na lista final da apresentação
\setbeamertemplate{bibliography item}{\insertbiblabel}


%----------------------------------------------------------------------------------------
%    TITLE PAGE
%----------------------------------------------------------------------------------------

\title{Vazamento}
%\subtitle{Licensed under CC-BY-SA}

\author{%
Thiago Abreu \and Evandro Rocha  \\%
Caio Lacoste \and Pedro Ivo %
}
\date{\today}
\institute{Laboratório de Processamento de Sinais (LPS)}


%----------------------------------------------------------------------------------------
%    PRESENTATION SLIDES
%----------------------------------------------------------------------------------------

\begin{document}
\begin{frame}[plain]
    \titlepage
\end{frame}

\begin{frame}{Conteúdo da Apresentação}
    % Throughout your presentation, if you choose to use \section{} and \subsection{} commands, these will automatically be printed on this slide as an overview of your presentation
    \tableofcontents
\end{frame}


%------------------------------------------------
\section{Definição de protótipo fidedigno}
%------------------------------------------------

\begin{frame}[plain]
    \vfill
    \centering
    {\Large\textcolor{dexpiBlue}{\textbf{Seção 1}}}\\[0.8em]
    {\LARGE\textbf{Definição de Protótipo Fidedigno}}\\[0.6em]
    \vfill
\end{frame}


%------------------------------------------------
\section{Tipos dos Gêmeos Digitais}
%------------------------------------------------

\begin{frame}[plain]
    \vfill
    \centering
    {\Large\textcolor{dexpiBlue}{\textbf{Seção 2}}}\\[0.8em]
    {\LARGE\textbf{Tipos dos Gêmeos Digitais}}\\[0.6em]
    \vfill
\end{frame}

\begin{frame}{Tipos dos Gêmeos Digitais}
    Algumas classificações de Gêmeos Digitais (GDs):
    \vspace{0.4em}
    \centering
    \footnotesize
    \begin{tabular}{p{3.5cm}p{3.0cm}p{5.3cm}}
        \toprule
        \textbf{Foco} & \textbf{Pioneiros / Referência} & \textbf{Conceito Central} \\
        \midrule
        1. Ciclo de Vida & Grieves \cite{grieves2017dt} & DTP, DTI e DTA \\
        2. Fluxo de Dados & Kritzinger \cite{fuller2020dt} & Modelo, Sombra e Gêmeo Digital \\
        3. Granularidade & Vohra \cite{vohra2022dt} & Componente, Ativo, Sistema, Processo \\
        4. Cognitivos (CDT) & El Adl \cite{eladl2016cdt} / Zheng \cite{zheng2022emergence} & Raciocínio, Grafos de Conhecimento \\
        5. Generativo (GenDT) & Xia \cite{xia2024aas} & Modelagem Generativa, Agentes (LLMs) \\
        \bottomrule
    \end{tabular}
\end{frame}

\begin{frame}{Gêmeo Digital Cognitivo (CDT)}
    Proposto conceitualmente por Ahmed El Adl \cite{eladl2016cdt} e formalizado por Zheng et al. \cite{zheng2022emergence}, o CDT estende o gêmeo tradicional adicionando capacidades de raciocínio, semântica e autoadaptação:
    \vspace{0.4em}
    \begin{columns}[T]
        \begin{column}{0.48\textwidth}
            \begin{block}{O Cérebro do CDT (Ontologias)}
                \footnotesize
                \begin{itemize}
                    \item \textbf{Grafos de Conhecimento:} Mapeiam o contexto e as inter-relações lógicas de todo o sistema físico.
                    \item \textbf{Raciocínio Lógico:} Dedução de causas e efeitos (ex: diagnosticar se uma anomalia é falha de leitura do sensor ou avaria real).
                \end{itemize}
            \end{block}
        \end{column}
        \begin{column}{0.48\textwidth}
            \begin{exampleblock}{Capacidades Cognitivas}
                \footnotesize
                Mapeamento da ciência cognitiva humana:
                \begin{itemize}
                    \item \textbf{Percepção e Atenção:} Filtragem de ruídos para foco em metas críticas.
                    \item \textbf{Memória e Aprendizado:} Atualização autônoma de modelos com base no desgaste temporal da máquina.
                \end{itemize}
            \end{exampleblock}
        \end{column}
    \end{columns}
\end{frame}

\begin{frame}{Casos de Sucesso do CDT em Operação}
    Aplicações práticas documentadas na literatura científica comprovam o valor dos CDTs:
    \vspace{0.4em}
    \begin{itemize}\setlength\itemsep{4pt}
        \item \textbf{Manufatura de Tubos de Aço \cite{zheng2022emergence}:} Um CDT de monitoramento mecânico utilizou aprendizado autônomo de vida útil dos componentes, gerando uma \textbf{redução de 10\% no consumo de energia} e minimizando paradas.
        \item \textbf{Montagem de Transmissões (PTU) \cite{zheng2022emergence}:} Integração de CDT com sensores heurísticos reduziu refugo industrial, elevando a precisão da detecção de falhas de montagem em \textbf{9\%}.
        \item \textbf{Prédios Inteligentes e Climatização (HVAC) \cite{elmokhtari2022hvac}:} Um CDT de controle preditivo realizou previsões de carga térmica 24 horas antes e ajustou dinamicamente ar-condicionado e luzes com base em presença, cortando desperdícios.
    \end{itemize}
\end{frame}

\begin{frame}{O Gêmeo Digital Generativo (GenDT)}
    A integração de IA Generativa (LLMs e modelos de difusão) foca em automatizar a criação do gêmeo e facilitar a interação \cite{xia2024aas}:
    \vspace{0.4em}
    \begin{itemize}\setlength\itemsep{5pt}
        \item \textbf{Modelagem Semântica Automatizada:} Agentes de IA leem especificações técnicas e manuais em PDF não estruturados e geram cascas administrativas padronizadas (Asset Administration Shell - AAS) automaticamente.
        \item \textbf{Twin Talk (Interface Conversacional):} O operador pergunta à máquina em linguagem natural (*``Qual o risco de quebra se rodarmos na velocidade máxima por 4 horas?''*). O LLM traduz isso em parâmetros de simulação, executa e retorna a resposta formatada.
    \end{itemize}
\end{frame}


\begin{frame}{Sinergia: Gêmeo Cognitivo + Generativo}
    A união dos dois modelos cria uma infraestrutura completa e robusta (sem alucinações de dados) \cite{mateev2023cdt}:
    \vspace{0.4em}
    \begin{columns}[T]
        \begin{column}{0.48\textwidth}
            \begin{block}{O Cérebro Racional (CDT \cite{zheng2022emergence})}
                \footnotesize
                \begin{itemize}
                    \item Baseado em \textbf{Grafos de Conhecimento}.
                    \item Garante que a IA respeite as regras físicas e lógicas exatas do mundo real.
                    \item Sem alucinações matemáticas ou operacionais.
                \end{itemize}
            \end{block}
        \end{column}
        \begin{column}{0.48\textwidth}
            \begin{exampleblock}{A Interface Conversacional (GenAI \cite{xia2024aas})}
                \footnotesize
                \begin{itemize}
                    \item Baseada em \textbf{LLMs e Agentes}.
                    \item Traduz comandos de linguagem natural em consultas formais ao Grafo de Conhecimento.
                    \item Facilita a interação de operadores humanos não técnicos.
                \end{itemize}
            \end{exampleblock}
        \end{column}
    \end{columns}
\end{frame}


%------------------------------------------------
\section{Maior aprofundamento nos Exemplos}
%------------------------------------------------

\begin{frame}[plain]
    \vfill
    \centering
    {\Large\textcolor{dexpiBlue}{\textbf{Seção 3}}}\\[0.8em]
    {\LARGE\textbf{Maior Aprofundamento nos Exemplos}}\\[0.6em]
    \vfill
\end{frame}



%------------------------------------------------
\section{Novos Exemplos}
%------------------------------------------------

\begin{frame}[plain]
    \vfill
    \centering
    {\Large\textcolor{dexpiBlue}{\textbf{Seção 4}}}\\[0.8em]
    {\LARGE\textbf{Novos Exemplos}}\\[0.6em]
    \vfill
\end{frame}


%------------------------------------------------
\section{Referências}
%------------------------------------------------

\begin{frame}[allowframebreaks]
\frametitle{Referências}
    \footnotesize
    \bibliographystyle{ieeetr}
    \bibliography{reference.bib}
    
\end{frame}

\begin{frame}
    \Huge{\centerline{\textbf{Chegamos ao Fim}}}
    \Huge{\centerline{\textbf{Obrigado pela sua Atenção}}}
\end{frame}


\end{document}